\documentclass{article}
\usepackage{iclr2026_conference,times}
\input{math_commands.tex}
\usepackage{booktabs}
\usepackage{graphicx}
\usepackage{hyperref}
\usepackage{url}

\title{Where Bits Matter in DINO-WM Planning:\\A Consolidated ICLR-Style Report from Completed M3 Experiments}
\author{}

\begin{document}
\maketitle
\lhead{Under review as a workshop paper at ICLR 2026 (ES-Reasoning)}

\begin{abstract}
This paper consolidates all completed quantization experiments currently available in our M3 pipeline for DINO-WM planning on Wall.
The central question is whether low-bit planning quality depends only on bitwidth, or also on where precision is allocated.
We report results from a strict run with 22 completed cells (66 evaluated episodes): a 20-cell core grid over FP16, uniform/mixed INT8, and uniform/mixed INT4 across two budgets and two seeds, plus a two-cell encoder localization block (tail vs head FP16 retention at 4-bit).
We also include one earlier pilot FP16 artifact as context.
The evidence is clear that 8-bit variants are robust and 4-bit variants are fragile.
However, mixed INT4 is not uniformly superior to uniform INT4 across both budgets in this completed run; its advantage is budget-dependent.
This paper provides the finalized tables/plots for the completed data, explains each visualization, and identifies the next experiments needed to convert budget-dependent signals into a stronger generalized claim.
\end{abstract}

\section{Introduction}
Recent world-model planners derive control from latent dynamics and representation quality.
At the same time, low-bit quantization is increasingly used for deployment under strict memory and compute constraints.
Our working hypothesis is that at low precision, module-aware allocation matters: preserving precision in representation-critical components may outperform uniform quantization at the same nominal bitwidth.

This report focuses on completed experiments only.
It is not a speculative plan.
Every number and figure is generated from saved outputs in the run directories.

\section{Experimental Protocol and Data Consolidation}
\paragraph{Task and model.}
We evaluate DINO-WM on Wall using paired-goal planning episodes.

\paragraph{Hardware and backend.}
All consolidated runs here use Apple M3 with MPS backend.

\paragraph{Completed runs included in this paper.}
\begin{table}[t]
\centering
\small
\begin{tabular}{lccc}
\toprule
Artifact & Cells & Episodes/Cell & Total Episodes \\
\midrule
Core grid (5 variants $\times$ 2 budgets $\times$ 2 seeds) & 20 & 3 & 60 \\
Where-bits block (tail50 vs head50, bA seed1) & 2 & 3 & 6 \\
Prior pilot FP16 context (m3\_main2exp\_v2) & 1 & 10 & 10 \\
\bottomrule
\end{tabular}
\caption{Inventory of completed artifacts consolidated in this paper.}
\label{tab:inventory}
\end{table}

Table~\ref{tab:inventory} clarifies what is and is not pooled.
The strict run contributes 66 episodes total.
The pilot FP16 row is reported separately as context because its protocol differs (10 episodes vs 3 per strict cell).

\section{Main Results}
\subsection{Overall Regime Structure: 8-bit Stable, 4-bit Fragile}
\begin{figure}[t]
\centering
\includegraphics[width=0.98\linewidth]{../figures_m3_main2_expansion/ m3_main2exp_v3/interim_metrics/interim_frontier_size_vs_success.pdf}
\caption{Accuracy-memory frontier from completed strict-run outputs. Marker color encodes variant family and annotations include budget ID.}
\label{fig:frontier}
\end{figure}

\begin{table}[t]
\centering
\small
\begin{tabular}{lccc}
\toprule
Variant (strict run aggregate) & Mean Success & Mean Time (s) & Mean Size (MB) \\
\midrule
FP16 & 0.8333 & 278.45 & 204.99 \\
Uniform INT8 & 0.7500 & 301.33 & 87.72 \\
Mixed INT8 & 0.8333 & 295.90 & 148.31 \\
Uniform INT4 & 0.1667 & 294.38 & 68.12 \\
Mixed INT4 & 0.1667 & 306.91 & 138.84 \\
\bottomrule
\end{tabular}
\caption{Core strict-run aggregate over bA/bB and seeds 1/2.}
\label{tab:core}
\end{table}

Figure~\ref{fig:frontier} and Table~\ref{tab:core} together show the first major pattern.
The 8-bit variants remain strong and close to FP16 in success, while both 4-bit variants drop sharply.
The frontier plot is useful because it shows this is not only an accuracy story; it also exposes memory compression tradeoffs directly.
In this run, the smallest model (uniform INT4) is also the weakest among the shown families.

\subsection{Budget-Conditioned Behavior at 4-bit}
\begin{figure}[t]
\centering
\includegraphics[width=0.98\linewidth]{../figures_m3_main2_expansion/ m3_main2exp_v3/interim_metrics/interim_budget_variant_bars.pdf}
\caption{Budget-wise success bars with seed-level points for key variants.}
\label{fig:budgetbars}
\end{figure}

\begin{figure}[t]
\centering
\includegraphics[width=0.88\linewidth]{../figures_m3_main2_expansion/ m3_main2exp_v3/interim_metrics/interim_mixed_vs_uniform_int4_delta.pdf}
\caption{Budget-level delta $\Delta$success = mixed INT4 $-$ uniform INT4 from strict-run means.}
\label{fig:deltabar}
\end{figure}

\begin{table}[t]
\centering
\small
\begin{tabular}{lccc}
\toprule
Budget & Uniform INT4 & Mixed INT4 & $\Delta$ (Mixed $-$ Uniform) \\
\midrule
bA & 0.1667 & 0.3333 & +0.1667 \\
bB & 0.1667 & 0.0000 & -0.1667 \\
\bottomrule
\end{tabular}
\caption{Budget-conditioned 4-bit comparison over seeds 1 and 2.}
\label{tab:budget}
\end{table}

Figure~\ref{fig:budgetbars} is the most important diagnostic for the core claim.
It shows that mixed INT4 is better than uniform INT4 on bA, but worse on bB.
Figure~\ref{fig:deltabar} and Table~\ref{tab:budget} make this explicit as signed deltas.
So the completed data does \emph{not} support a universal statement that mixed INT4 dominates uniform INT4 in this setup.
Instead, the effect is budget-conditioned.

\subsection{Paired Statistics and Where-Bits Localization}
\begin{figure}[t]
\centering
\includegraphics[width=0.98\linewidth]{../figures_m3_main2_expansion/ m3_main2exp_v3/interim_metrics/interim_paired_delta_ci.pdf}
\caption{Paired deltas with 95\% CIs from episode-level paired analysis JSON outputs.}
\label{fig:pairedci}
\end{figure}

\begin{figure}[t]
\centering
\includegraphics[width=0.92\linewidth]{../figures_m3_main2_expansion/ m3_main2exp_v3/interim_metrics/interim_success_heatmap.pdf}
\caption{Variant-by-budget success heatmap, including where-bits variants (tail50 and head50).}
\label{fig:heatmap}
\end{figure}

Paired statistics refine the mean-level view.
Figure~\ref{fig:pairedci} shows:
(1) bA mixed-vs-uniform INT4 delta is positive with CI touching zero,
(2) bB delta is negative with CI touching zero,
(3) tail50-vs-head50 delta is exactly zero in this sample.
The confidence intervals are wide because paired units are small (6 for mixed-vs-uniform per budget, 3 for tail-vs-head).

Figure~\ref{fig:heatmap} helps interpret the full matrix quickly.
It confirms strong 8-bit performance, low 4-bit performance, and near-identical tail/head localization scores in the one-cell where-bits block (both 0.6667 on bA seed1).
This localization result should be treated as provisional because it comes from a single seed and single budget.

\section{What This Consolidated Evidence Supports}
From completed outputs, we can support the following claims:
\begin{itemize}
  \item 8-bit quantization (uniform/mixed) remains close to FP16 in this Wall setup.
  \item Moving to 4-bit substantially reduces success in aggregate.
  \item Mixed INT4 vs uniform INT4 is budget-sensitive rather than consistently positive in this completed run.
  \item The current tail-vs-head localization check is inconclusive at present sample size.
\end{itemize}

\section{Limitations}
The strict run uses small per-cell episode count (3), one environment (Wall), one checkpoint family, and limited where-bits coverage.
These choices were intentional for runtime control, but they widen uncertainty.
As a result, this paper should be interpreted as a consolidated completed-run report, not a final generalized theorem about all planning regimes.

\section{Conclusion}
This ICLR-style consolidated report combines all completed artifacts to date and presents a consistent story:
bitwidth alone is not sufficient to explain behavior near the low-bit frontier, but current evidence is not yet enough to claim universal mixed-INT4 superiority.
The data instead points to budget-conditioned effects, motivating the next phase: higher episode counts and broader budget sweeps to determine where mixed allocation reliably wins.

\end{document}
